\section*{Bab \ref{ch:g6:lines} --- Garis, Lingkaran dan Sudut}

\begin{solution}{exo:g6:lines:1}
Lukisan bebas. Ruas garisnya berhenti di $A$ dan $B$; sinar garisnya
bermula di $C$ lalu berlanjut melewati $A$; garisnya berlanjut ke
kedua arah melewati $B$ dan $C$.
\end{solution}

\begin{solution}{exo:g6:lines:2}
``$[AB] = [BA]$'': \emph{benar} --- bagian di antara kedua titik itu
tidak bergantung pada urutannya.

``$[AB) = [BA)$'': \emph{salah} --- $[AB)$ bermula di $A$, $[BA)$
bermula di $B$; keduanya menunjuk ke arah yang berlawanan.

``$(AB) = (BA)$'': \emph{benar} --- kedua nama itu menjelaskan garis
tak terbatas yang sama.
\end{solution}

\begin{solution}{exo:g6:lines:3}
Langkahnya: geser penggaris siku sepanjang $d$ sampai tepi tegak
lurusnya melalui $P$; gambarlah \omterm{def:g6:lines:perp}{garis tegak lurus} itu, sebut saja $p$.
Lalu gambarlah \omterm{def:g6:lines:perp}{garis tegak lurus} terhadap $p$ melalui $P$ dengan cara
yang sama: menurut \cref{prop:g6:lines:facts} garis itu \omterm{def:g6:lines:perp}{sejajar} dengan $d$.
\end{solution}

\begin{solution}{exo:g6:lines:4}
$a$ dan $b$ sama-sama \omterm{def:g6:lines:perp}{tegak lurus} terhadap garis yang sama $c$, jadi
$a \parallel b$ (fakta 1). Untuk $e$ dan $c$: keduanya \omterm{def:g6:lines:perp}{tegak lurus}
terhadap garis yang sama $a$ (diketahui $e \perp a$, dan $a \perp c$
berarti $c \perp a$ juga), jadi fakta 1 berlaku lagi: $e \parallel c$.
\end{solution}

\begin{solution}{exo:g6:lines:5}
$OM = 3$ cm tepat ($M$ pada lingkaran); $ON < 3$ cm ($N$ di dalam);
$OP > 3$ cm ($P$ di luar).
\end{solution}

\begin{solution}{exo:g6:lines:6}
\omterm{def:g3:shapes:shapes}{Jari-jari} $= 9 \div 2 = 4.5$ cm. \quad \omterm{def:g4:geometry:circle}{Diameter} $= 2 \times 2.6 = 5.2$ cm.
\end{solution}

\begin{solution}{exo:g6:lines:7}
$\widehat{SRT} = 72^\circ$: lancip. \quad
$\widehat{ABC} = 148^\circ$: tumpul. \quad
$\widehat{MON} = 90^\circ$: siku-siku. (\omterm{def:g5:solids:def}{Titik sudutnya} selalu huruf
yang di tengah.)
\end{solution}

\begin{solution}{exo:g6:lines:8}
Ukurannya bergantung pada segitiga yang digambar, tetapi \omterm{def:g1:addition:def}{jumlah}
ketiga sudutnya selalu (sangat dekat dengan) $180^\circ$ --- \omterm{ex:g1:subtraction:difference}{selisih}
yang kecil berasal dari ketidaktelitian pengukuran.
\Cref{ch:g7:triangles} membuktikan bahwa \omterm{def:g1:addition:def}{jumlahnya} tepat $180^\circ$.
\end{solution}

\begin{solution}{exo:g6:lines:9}
Lukisan bebas. Untuk memperoleh $130^\circ$ dari $65^\circ$ tanpa
busur derajat: salinlah sudut $65^\circ$ itu dua kali berdampingan
($65 + 65 = 130$), misalnya dengan kertas kalkir atau dengan menyalin
sudut memakai jangka dan penggaris.
\end{solution}

\begin{solution}{exo:g6:lines:10}
Titik yang berjarak $3$ cm dari $A$ membentuk lingkaran berpusat $A$
dengan \omterm{def:g3:shapes:shapes}{jari-jari} $3$ cm; titik yang berjarak $2$ cm dari $B$ membentuk
lingkaran berpusat $B$ dengan \omterm{def:g3:shapes:shapes}{jari-jari} $2$ cm. Titik yang diminta
adalah perpotongan kedua lingkaran itu: dua titik jika keduanya
berpotongan ($AB$ tepat di antara $1$ dan $5$ cm), satu titik jika
keduanya bersinggungan ($AB = 5$ cm atau $AB = 1$ cm), tidak ada jika
keduanya terlalu berjauhan atau yang satu di dalam yang lain.
\end{solution}

\begin{solution}{exo:g6:lines:11}
Ada $12$ tanda jam yang membelah satu putaran penuh menjadi $12$
sudut sebesar $360 \div 12 = 30^\circ$. Pada pukul 3 kedua jarumnya
merentang $3$ tanda: $90^\circ$ (\omterm{def:g3:shapes:rightangle}{sudut siku-siku}). Pada pukul 5:
$5 \times 30 = 150^\circ$. Pada pukul 6: $180^\circ$ (sudut lurus).
\end{solution}

\begin{solution}{exo:g6:lines:12}
Gambarlah dua busur berjari-jari $6$ cm yang berpusat di $A$ dan di
$B$; titik potongnya adalah $C$. Ketiga sisinya berukuran $6$ cm: jadi
segitiganya sama sisi, dan setiap sudutnya berukuran $60^\circ$
(pengukurannya membenarkan hal itu). Dugaan: segitiga sama sisi punya
tiga sudut yang sama besar, masing-masing $60^\circ$.
\end{solution}

\begin{solution}{pb:g6:lines:1}
\textbf{1.} Setengah putaran: $360 \div 2 = 180^\circ$. \omterm{def:g3:division:half}{Seperempat}:
$360 \div 4 = 90^\circ$. Seperdua belas: $360 \div 12 = 30^\circ$.

\textbf{2.} Dua belas sudut sama besar yang mengisi $360^\circ$:
$30^\circ$ di antara tanda yang bertetangga. Pada pukul 3 kedua
jarumnya merentang $3$ tanda: $3 \times 30 = 90^\circ$ --- \omterm{def:g3:shapes:rightangle}{sudut
siku-siku}, seperti yang ditemukan pada \cref{exo:g6:lines:11}.

\textbf{3.} Pada pukul 1: $1 \times 30 = 30^\circ$. Pada pukul 4:
$4 \times 30 = 120^\circ$. Pada pukul 7 kedua jarumnya merentang $7$
tanda di satu sisi, $7 \times 30 = 210^\circ$, jadi sudut di antara
kedua jarumnya adalah sisi yang lain: $360 - 210 = 150^\circ$.

\textbf{4.} $360^\circ$ dalam $60$ menit: $360 \div 60 =
6^\circ$ setiap menit.

\textbf{5.} $30^\circ$ dalam $60$ menit: $30 \div 60 = 0.5^\circ$
setiap menit --- setengah derajat. Jarum jam merayap, jarum menit
melangkah lebar.

\textbf{6.} Dalam $30$ menit sejak pukul 3:00, jarum jam sudah
berpindah $30 \times 0.5 = 15^\circ$ melewati angka $3$: ia duduk
tepat di tengah antara angka $3$ dan $4$, pada $90 + 15 = 105^\circ$
dari angka $12$. Jarum menit menunjuk angka $6$: $180^\circ$. Sudut di
antara kedua jarumnya: $180 - 105 = 75^\circ$.

\textbf{7.} Perhitungan mengalahkan terkaan: jarum menit ada di
$180^\circ$, tetapi jarum jam sudah meninggalkan angka $6$ --- ia ada
di $6 \times 30 + 30 \times 0.5 = 180 + 15 = 195^\circ$. Sudut di
antara kedua jarumnya $195 - 180 = 15^\circ$: dekat, tetapi tidak
bertumpuk.

\textbf{8.} Jarum menit di angka $3$: $90^\circ$. Jarum jam:
$9 \times 30 + 15 \times 0.5 = 270 + 7.5 = 277.5^\circ$. \omterm{ex:g1:subtraction:difference}{Selisihnya}:
$277.5 - 90 = 187.5^\circ$, jadi sudut di antara kedua jarumnya
$360 - 187.5 = 172.5^\circ$ --- hampir sudut lurus.

\textbf{9.} Pada pukul 1:00 jarum menit ($0^\circ$) berada
\emph{di belakang} jarum jam ($30^\circ$). Pada pukul 1:10 jarum
menit ($60^\circ$) berada \emph{di depannya}
($30 + 10 \times 0.5 = 35^\circ$). Jarum menit berjalan lebih cepat
($6^\circ$ setiap menit melawan $0.5^\circ$), jadi antara pukul 1:00
dan 1:10 ia menyusul lalu melewati jarum jam --- pada saat melewati
itulah kedua jarumnya tepat bertumpuk.

\textbf{10.} Penyusulan yang sama terjadi berulang-ulang: kedua
jarumnya bertemu tepat pada pukul 12:00, lalu kira-kira pada 1:05,
2:11, 3:16, 4:22, 5:27, 6:33, 7:38, 8:44, 9:49 dan 10:55. Pertemuan
yang mungkin diharapkan antara pukul 11 dan 12 jatuh tepat \emph{pada}
pukul 12:00, yang sudah memulai putaran berikutnya --- jadi dalam dua
belas jam kedua jarumnya berimpit $11$ kali, bukan $12$. Dalam satu
hari penuh: $22$ kali.

\textbf{11.} Sudut yang berbagi satu \omterm{def:g5:solids:def}{titik sudut}, tanpa tumpang tindih
dan tanpa celah, memenuhi satu putaran penuh di sekeliling titik itu:
menyapu sekali mengelilingi titik itu melewati setiap sudutnya tepat
sekali, dan satu sapuan penuh adalah $360^\circ$. Jadi ukurannya
berjumlah $360^\circ$ --- berapa pun banyak sudutnya dan berapa pun besarnya.

\textbf{12.} Setiap jam dalam sehari bernilai
$360 \div 24 = 15^\circ$. Jadi:
\[
\text{tidur } 9 \times 15 = 135^\circ, \quad
\text{sekolah } 90^\circ, \quad
\text{bermain } 45^\circ, \quad
\text{makan } 30^\circ, \quad
\text{lainnya } 60^\circ,
\]
dan $135 + 90 + 45 + 30 + 60 = 360^\circ$: lingkarannya penuh, tanpa
celah dan tanpa tumpang tindih (pertanyaan 11).

\textbf{13.} Gambarlah sebuah lingkaran dengan jangka beserta satu
\omterm{def:g3:shapes:shapes}{jari-jari} (garis awalnya). Letakkan pusat busur derajat pada pusat
lingkaran itu, garis $0^\circ$-nya pada \omterm{def:g3:shapes:shapes}{jari-jari} tadi, lalu tandai
$135^\circ$; gambarlah \omterm{def:g3:shapes:shapes}{jari-jari} yang baru: juring tidur selesai.
Lalu letakkan garis $0^\circ$-nya pada \emph{\omterm{def:g3:shapes:shapes}{jari-jari} itu} lalu
tandai $90^\circ$ untuk sekolah, dan seterusnya --- setiap juring
bermula di tempat juring sebelumnya berakhir. Setelah juring terakhir,
gambarnya menutup tepat pada \omterm{def:g3:shapes:shapes}{jari-jari} awalnya.

\textbf{14.} $90^\circ$ adalah $\frac{90}{360} = \frac14$ cakramnya
(\cref{met:g6:fractions:simplify}), jadi ia mewakili \omterm{def:g3:division:half}{seperempat} hari:
$24 \div 4 = 6$ jam.

\textbf{15.} Jarum jam menempuh satu putaran dalam $12$ jam: $2$
putaran sehari. Jarum menit, satu putaran setiap jam: $24$ putaran.
Jarum detik, satu putaran setiap menit: $24 \times 60 = 1\,440$
putaran sehari.

\end{solution}
